% \begin{savequote}
% \includegraphics[width=7cm]{./images/gabriel-sebastian-aass/St-Anna-Kinderspital_IMG_2000_00800.jpg}
% \end{savequote}
\def\ChapterCode{KIN}
\chapter
[\CO{[KIN]} \emph{Spezielle Patientengruppe:} Kinder]
{\emph{Spezielle Patientengruppe:} Kinder}
\label{chp:KIN}

\ChapterIntro
{%
~~~
}
{%
\begin{description}
  \item[Maintainer:] Sebastian Gabriel
  \item[Autoren:] Diverse
  \item[Reviewer:] Standard-Reviewprozess
  \item[Version:] {\VarDocVersionTypeName} ({\VarDocVersionTypeDescription})
  \item[SHA1:] {\DocGitCommit}
\end{description}

{\TxTeilkapitelAASS}
}
%%%%%%%%%%%%%%%%%%%%%%%%%%%%%%%%%%%%%%%%%%%%%%%%%%%%%%%%%%%
%%% Pädiatrie                                           %%%
%%%%%%%%%%%%%%%%%%%%%%%%%%%%%%%%%%%%%%%%%%%%%%%%%%%%%%%%%%%

\section{\HeadingKeyword{Umgang} mit Kindern}
\Text
{{\TxBeschreibung}}%1 Überschrift
{%
Der Umgang soll prinzipiell schonend erfolgen, Aufregung
soll nach Möglichkeit vermieden werden. Wichtig ist, dem
Kind eine Bezugsperson zu geben; Die Vertrauensperson bzw.
Eltern müssen unbebedingt in die (Be-)Handlung einbezogen
werden. Nach Möglichkeit soll das Kind während eng bei
Vertrauensperson belassen werden (auf den Schoss setzen,
\ldots). 
\nocite{LPN3}
}%2 Text



\section{Anatomische und physiologische  \HeadingKeyword{Besonderheiten}}
\Text
{{\TxBeschreibung}}%1 Überschrift
{%
\label{topic:kind-besonderheiten}%
Kinder unterscheiden sich wesentlich von erwachsenen
Patienten, nicht nur in ihrer Fähigkeit ihre Umwelt und
gesetzte Handlungen zu verstehen, sondern auch hinsichtlich
ihrer Anatomie und Physiologie.
\FullPrettyRef{tab:altersgruppen-kinder}, zeigt
eine Übersicht über die Einteilung von Kindern nach
Altersgruppen.

}%2 Text

\begin{Synopsis}
  \SynopsisItem{Kinder sind keine kleinen Erwachsenen!}
\end{Synopsis}



\Table
{Übersicht: Altersgruppen}%1 Titel
{In der Literatur kommt es
gelegentlich zu abweichenden Definitionen. So ist z.\,B. in den
ERC-Richtlinien zur Reanimation ein Neugeborenes nur bis
unmittelbar nach der Geburt als solches
klassifiziert.}%2 Beschreibung
{\label{tab:altersgruppen-kinder}}%3 Label
{}%4 Quelle
{}%5 Lizenz
{%
\begin{tabulary}{\linewidth}{LLL}\addlinespace\toprule\addlinespace
\noindent\TabHeading{Bezeichnung}&\noindent\TabHeading{Alter}&\noindent\TabHeading{alternativ}
\\\addlinespace\midrule\addlinespace
\noindent\TabSub{Neugeborenes} & bis zum 28. Lebenstag & unmittelbar nach der Geburt
\\\addlinespace
\noindent\TabSub{Säugling} & bis zum Ende des 1. Lebensjahr
\\\addlinespace
\noindent\TabSub{Kleinkind} & \VonBis{1}{5} Jahre
\\\addlinespace
\noindent\TabSub{Schulkind} & \VonBis{6}{13} Jahre
\\\addlinespace
\noindent\TabSub{Jugendlicher} & \VonBis{14}{18} Jahre
\\\bottomrule
\end{tabulary}
}%6 Tabelle
{}%7 Float

% \begin{PictureSeriesWide}{}{Pädiatrisches Notfalllineal}
% \PictureSeriesRow{2}
% {./images/ASBOE-Grassl-gjh-1/CRW_0409-00441pt}
% {Ein pädiatrisches Notfalllineal  \SourcePicture{ASBÖ/gjh}}
% {./images/ASBOE-Grassl-gjh-1/CRW_0412-00441pt}
% {Mittels eines pädiatrischen Notfalllineals können größenspezifische Normalwerte und Dosierungen von Medikamenten abgelesen werden. \SourcePicture{ASBÖ/gjh}}
% {}
% {}
% {}
% {}
% \end{PictureSeriesWide}

% TODO Grassl Foto: Pädiatrisches Notfalllineal



\Table
{Anatomische und physiologische Besonderheiten von Kindern}%1 Titel
{}%2 Beschreibung
{}%3 Label
{}%4 Quelle
{}%5 Lizenz
{%
\begin{tabulary}{\linewidth}{LL}\toprule\addlinespace
\noindent\TabSub{Atemfrequenz}
&
Neugeborenes: \hfill \VonBis{30}{50}\PerMin*

Säugling: \hfill \VonBis{20}{30}\PerMin*

Kleinkind: \hfill \VonBis{20}{30}\PerMin*

Beatmung bis zum Säugling in  Neutralposition (Kopf nicht
überstrecken)
\\\addlinespace\midrule\addlinespace
\noindent\TabSub{Zeichen der Atemnot}%
\label{topic:atemnot-zeichen-kinder}%
&
Einziehungen zwischen Rippen

\T{Nasenflügeln}
beschleunigte Atmung,
Atemgeräusche
\\\addlinespace\midrule\addlinespace
\noindent\TabSub{Pulstasten beim Säugling}
&
Oberarmarterie (Oberarm Innenseite), Achsel
\\\addlinespace\midrule\addlinespace
\noindent\TabSub{Herzfrequenz}
&
Neugeborenes: \hfill \VonBis{140}{180}\PerMin*\newline\hfill(Lebensgefahr unter 80\PerMin*)

Säugling: \hfill \VonBis{110}{160}\PerMin*\nocite{Lissauer2007}

Kleinkind: \hfill \VonBis{95}{140}\PerMin*\nocite{Lissauer2007}

Schulkind: \hfill \VonBis{80}{120}\PerMin*\nocite{Lissauer2007}

Bradykardie  ist fast immer  hypoxiebedingt 
\\\addlinespace\midrule\addlinespace
\noindent\TabSub{Systolischer Blutdruck}
&
Neugeborenes: \hfill 75\MmHg*\nocite{Soehnke200607}\nocite{Lissauer2007}

Säugling: \hfill \VonBis{80}{90}\MmHg*\nocite{Soehnke200607}

Kleinkind: \hfill 95\MmHg*\nocite{Lissauer2007}

Schulkind:
\hfill \VonBis{100}{110}\MmHg*\nocite{Soehnke200607}\nocite{Lissauer2007}

Jugendlicher: \hfill 120\MmHg*\nocite{Soehnke200607}
\\\addlinespace\midrule\addlinespace
\noindent\TabSub{Wasser-, Elektrolyt- und Wärmehaushalt}
&
Kinder haben im Verhältnis zu ihrem  Körpergewicht eine
größere Körperoberfläche  als Erwachsene!

höherer Wasseranteil

Austrocknung und Unterkühlung schneller  möglich!
\\\addlinespace\bottomrule
\end{tabulary}
}%6 Tabelle
{}%7 Float

% TODO Grassl Foto: Diagnostik, Therapie bei Kindern (RR, Temperatur, HWS-Schiene)

\subsection{Allgemeines zu \HeadingKeyword{Erkrankungen im Kindes- und
Jugendalter}}

\Text
{{\TxBeschreibung}}%1 Überschrift
{%
Grundsätzlich können Kinder viele der bereits erwähnten
Erkrankungen auch bekommen. Bei Kindernotfällen können als
zusätzlicher Stressfaktor für den Rettungsdienst stark
emotional betroffene Angehörige sein. Daher gilt es
besonders ruhig und gezielt zu handeln. Sie sollten sich der
Verantwortung bewusst sein, es geht um die Rettung eines
jungen Menschen mit langer Lebenserwartung.
}%2 Text



\section{Speziell die Kindheit und Jugend betreffende \HeadingKeyword{Erkrankungen}}




\subsection{\HeadingKeyword{Fremdkörperaspiration}}


\Text
{{\TxBeschreibung}}%1 Überschrift
{%
Die Fremdkörperaspiration und mechanische Atemwegsverlegung wird unter
\FullPrettyRef{topic:mechanische-atemwegsverlegung},
behandelt.
}%2 Text


\subsection{Akute obstruktive \HeadingKeyword{Laryngitis}, Pseudokrupp}
\label{topic:laryngitis}
\label{topic:pseudokrupp}

% TODO J05.- Akute obstruktive Laryngitis [Krupp] und Epiglottitis überarbeiten

\TextSlide{\TxBeschreibung}
{%
Bei der Laryngitis
(Pseudokrupp\index{Pseudokrupp}\footnote{\enquote{Pseudokrupp}
bezieht sich auf die Heiserkeit (griech.: pseudo für
„unecht“; schottisch: croup,
„Heiserkeit“\cite{wiki:pseudokrupp}). Der echte Krupp
bezeichnet die Kehlkopfentzündung bei Diphtherie.}) handelt
es sich um eine Entzündung des Kehlkopfes. Es können auch die
oberen Atemwege durch das Anschwellen der Schleimhäute
behindert sein (Pseudokrupp/Kruppsyndrom). Diese Infektion
wird durch Viren verursacht. 
}{
\begin{itemize}
  \item Entzündung des Kehlkopfes
\end{itemize}
}{}{}{}


\TextSlide{\TxSymptome}
{%
Das Leitsymptom ist ein typischer, \EE{trockener, bellender
Husten}. Oft ist auch ein pfeifendes Geräusch beim Einatmen
\E{(inspiratorischer Stridor)} wahrnehmbar, und es besteht Heiserkeit und Dyspnoe. Fieber ist eher
diskret und steht nicht im Vordergrund. Die Laryngitis tritt
vor allem im Alter zwischen 3 Monaten und 6 Jahren auf, gut geheizte und trockene Räume fördern das Auftreten.

%\SlideCompensation{2}

}{
\begin{itemize}
    \item Trockener, bellender Husten
    \item Pfeifendes Geräusch beim Einatmen, Heiserkeit
    \item Dyspnoe
    \item Kaum Fieber
    \item Typisch:
    \begin{itemize}
      \item Alter: 3 Monate -- 6 Jahre
      \item warme, trockene Räume
    \end{itemize}
\end{itemize}
}{}{}{}

\begin{Danger}
  \item Bei Manipulation im Mund-/Rachenraum kann es zu einem Anschwellen und einer Verlegung der Atemwege kommen!
\end{Danger}

% M %%%%%%%%%%%%%%%%%%%%%%%%%%%%%%%%%%%%%%%%%%%%%%%%%%% M %
% Akute obstruktive Laryngitis
\ProcedurePrintDefault{KJ05000C}


\subsection{Akute \HeadingKeyword{Epiglottitis}}
\label{topic:epiglottitis}

\Text{\TxBeschreibung}
{%
Bei der Epiglottitis handelt es sich um eine \E{bakterielle
Entzündung} und daraus folgende \EE{Schwellung des
Kehldeckels}. Sie kommt eher selten vor, kann jedoch
\Ca{lebensbedrohlich} werden. Die Epiglottitis ist eine schwere
Erkrankung, mit richtiger  Behandlung besteht aber meist
eine gute Prognose.
}{%
\begin{itemize}
  \item Entzündung mit Schwellung des Kehldeckels
\end{itemize}
}{}{}{}



\TextSlide{\TxSymptome}
{%
Die Epiglottitis geht mit hohem Fieber einher. Es ist ein
Atemgeräusch bei der Einatmung wahrnehmbar
(\E{inspiratorischer Stridor}). Der Patient hält den Mund
offen und es ist ein Speichelfluss bemerkbar. Der Patient
ist eher ruhig und spricht kaum, er hat Angst.
}{%
\begin{itemize}
    \item Hoch fieberhafte Erkrankung
    \item Inspiratorischer Stridor
    \item Mund offen, Speichelfluss, spricht kaum!
    \item Angst!
\end{itemize}
}{}{}{}

\begin{Danger}
  \item Bei Manipulation im Mund-/Rachenraum kann es zu einem Anschwellen und einer Verlegung der Atemwege kommen!
\end{Danger}

% M %%%%%%%%%%%%%%%%%%%%%%%%%%%%%%%%%%%%%%%%%%%%%%%%%%% M %
% Akute Epiglottitis
\ProcedurePrintDefault{KJ05010C}


\subsubsection{\HeadingKeyword{Vergleich} Laryngitis vs. Epiglottitis}
 
Die Unterscheidung beider Krankheiten kann auch für den
erfahrenen Sanitäter schwer sein. Darum hier eine
Gegenüberstellung der Differenzialdiagnosen Laryngitis
und Epiglottitis:


\Table
{Gegenüberstellung Laryngitis vs. Epiglottitis}%1 Titel
{}%2 Beschreibung
{}%3 Label
{}%4 Quelle
{}%5 Lizenz
{%
\begin{tabulary}{\linewidth}{LL}
\toprule\addlinespace
\noindent\TabHeading{Laryngitis}
&
\noindent\TabHeading{Epiglottitis}
\\\addlinespace\midrule\addlinespace
Wenig bedrohlich
&
Lebensgefährlich
\\\addlinespace
Kind ist weinerlich
&
Ruhiges Kind, apathisch
\\\addlinespace
Kann schlucken
&
Speichelfluss
\\\addlinespace
Kaum Fieber
&
Hohes Fieber
\\\addlinespace
Kind schreit; weint evtl.
&
Sehr leise, verfällt zusehends
\\\addlinespace\bottomrule
\end{tabulary}
}%6 Tabelle
{}%7 Float

 
 
\subsection{SIDS: Sudden Infant Death Syndrome (\HeadingKeyword{Plötzlicher Kindstod})}
\index{plötzlicher Kindstod}
\index{SIDS}
\index{Sudden Infant Death Syndrome}
\label{topic:sids}

\TextSlide{\TxBeschreibung}
{%
Der plötzliche Kindstod  ist eine plötzliche, einschneidende
Katastrophe im Familienleben. Aus \EE{unbekannter Ursache}
-- und aus oft völliger Gesundheit -- hört das Kind
plötzlich auf zu atmen und verstirbt.
Dieses Schicksal ereilt Kinder im ersten bis vierten
Lebensmonat, es ist eine Häufung in den Monaten Jänner bis
März zu beobachten.

}{
\begin{itemize}
    \item Plötzlich, meist im Schlaf
    \item Vorwiegend \VonBis{1.}{4.} Lebensmonat
    \item Gehäuft Jänner bis März
    \item Bei Frühgeburten gehäuft
\end{itemize}
}{}{}{}

\TextSlide{Ursachen}
{%
Es wird viel über die möglichen Ursachen spekuliert. \Z{Als
mögliche Ursache wird oft ein unreifes Atemzentrum und eine
zentrale Atemregulationsstörung angeführt} Letzten Endes ist
die genaue \EE{Ursache unbekannt}, eine verlässliche
Vorhersage, welche Kinder betroffen sein werden, gibt es
nicht. Es wurden lediglich Risikogruppen wie Frühgeburten
und Zwillingsgeschwister von betroffenen Kindern,
identifiziert. Auch scheint die Lagerung in Bauchlage das
Risiko zu erhöhen.

Es muss betont werden, dass eine verlässliche Vorhersage
nicht möglich ist und \EE{die Eltern keine Schuld trifft}!

}{
\begin{itemize}
    \item Ursache unbekannt
    \item Risikogruppen
    \item Zentrale Atemregulationsstörung \Z{(unreifes 
    Atemzentrum?)}
    \item Eltern trifft keine Schuld!
\end{itemize}
}{}{}{}


% M %%%%%%%%%%%%%%%%%%%%%%%%%%%%%%%%%%%%%%%%%%%%%%%%%%% M %
% SIDS
\ProcedurePrintDefault{KR95090C}


\subsection{\HeadingKeyword{Ertrinkungsunfall}}\index{Ertrinkungsunfall}

\Text
{\TxBeschreibung}
{%
\DefinitionInPara{Als \EE{Ertrinken} wird der Tod in den
ersten 24\,h nach einem Ertrinkungsunfall bezeichnet; ansonsten besteht ein
\index{Beinahe-Ertrinken}\T{Beinahe-Ertrinken}.}
Kleinkinder sind besonders gefährdet, da sie meist nicht
schwimmen können und sich oft schnell und unbemerkt der
Aufmerksamkeit der Aufsichtsperson entziehen. Besonders
ungesicherte Gartenteiche, Seen, Schwimmbecken, Badeanstalten
oder auch Wannenbäder stellen eine Gefahrenquelle im Alltag
dar.}
{
\begin{itemize}
  \item Ertrinken: Tod innerhalb 24\Hour* nach Ertrinkungsunfall
  \item Sonst: Beinahe-Ertrinken
\end{itemize}
}
{}
{}
{}

\begin{Synopsis}
  \SynopsisItem{Nobody is dead until warm and dead!}
\end{Synopsis}


% M %%%%%%%%%%%%%%%%%%%%%%%%%%%%%%%%%%%%%%%%%%%%%%%%%%% M %
% Ertrinkungsunfall
\ProcedurePrintDefault{MT75010C}





\subsection{\HeadingKeyword{Krampfanfälle} im Kindesalter}
\label{topic:krampfanfaelle-kinder}

Grundsätzlich gelten auch bei Kindern die Ausführungen im
Abschnitt \ref{topic:zerebrale-krampfanfaelle}. Eine
Besonderheit stellt jedoch der \T{Fieberkrampf} dar.


\subsubsection{\HeadingKeyword{Fieberkrampf}}
\label{topic:fieberkrampf}

\TextSlide{\TxBeschreibung}
{%
Kleinkinder \EE{unter 5~Jahren} können bei Fieber sogenannte
Fieberkrämpfe entwickeln. Sie entsprechen einem vom Gehirn
ausgehenden Krampfanfall und sind grundsätzlich genau so zu
behandeln.
Mitunter ist die Neigung des Kindes zu Fieberkrämpfen schon
bekannt und es wurden bereits spezielle Zäpfchen für den
Bedarfsfall verschrieben (und bei Eintreffen eventuell schon
von den Eltern gegeben).
}{
\begin{itemize}
  \item Kinder {\textless} 5 Jahre
  \item Hohes Fieber
  \item Wie generalisierter Krampfanfall
\end{itemize}
}{}{}{}


% M %%%%%%%%%%%%%%%%%%%%%%%%%%%%%%%%%%%%%%%%%%%%%%%%%%% M %
% Fieberkrampf im Kindesalter
\ProcedurePrintDefault{KR56001C}


\subsection{\HeadingKeyword{Kindesmisshandlung}}
% 
% \Definition{Absichtliche psychische oder physische Gewalt gegen Kinder}
% 
% 
% Kindesmisshandlung kommt in \EE{allen sozialen Schichten} vor!

\TextSlide{\TxQuerverweise}
{%
\begin{itemize}
  \item \CO{[KRI]} \FullPrettyRef{topic:kindesmisshandlung}
\end{itemize}
}{%
\begin{itemize}
  \item \CO{[KRI]} \prettyref{topic:kindesmisshandlung}
\end{itemize}
}{}{}{}




\begin{Literary}
\fullcite{Lissauer2007}

\fullcite{SitzmannPaediatrie2}

\fullcite{NetterPaediatrie1}

\fullcite{Erc2010Sec6}

Weiters: \cite{
% SiewertChirurgie8,
AMLS3,
% Erc2005DeSec5,
% Erc2005Sec5,
% Erc2005DeSec6,
% Erc2005DeSec3,
% Erc2005Sec3,
BoehmerTaschRett6,
Fauci2008,
LPN3,Boecker2,
LippertAnatomie7,
% Nilsson1997,
RedelsteinerHRNS2,
% ForMed2,
StriebelAn7,
% Zeiler2001,
KlinkePhysio3,
DRAn3,
ForthPharma8,
LuellmannPharma16,
% ThierbachInterhospitaltransfer1,
AnCompact3,
RueckerBildatlas1}
\end{Literary}



\ChapterBibliography
\ChapterEnd